\section{术语消化：本期关键词索引}

本期概念很多，且许多词不是标准教科书用法。这里集中整理，方便读者复盘。

\begin{longtable}{p{0.24\textwidth}p{0.34\textwidth}p{0.33\textwidth}}
\toprule
术语 & 一句话解释 & 在本期中的作用 \\
\midrule
ONE2X & 王冠创办的 AI 产品工作室 & 本期实践样本。 \\
Mideo & ONE2X 的 AI 视频生成/处理产品 & 用来验证视频生成和生成系统。 \\
System1 & 模型压缩后的快速、本能反应能力 & 产品经理要设计模型能力和评测。 \\
System2 & 围绕模型组织的推理、流程、context 和生成系统 & 应用公司构建壁垒的主战场。 \\
Context Engineering & 设计模型生成前可用的有效上下文 & 把业务 know-how 转成模型可用 token。 \\
DSL & 领域特定语言/表达体系 & 用来描述视频生产方法和问题结构。 \\
Recipe & 可复用生产方法，不是单点内容 & 从内容生产转向 IP 和信任的关键。 \\
产品内生数据 & 产品使用中产生、之前不存在的数据 & 应用创业机会所在。 \\
生成系统 & 按意图和 recipe 组织内容生产的系统 & 对推荐系统的平台权力形成挑战。 \\
推荐系统 & 从已有库存中匹配内容给用户 & 互联网平台分发权的基础。 \\
分销平台 & 掌握内容分配和流量分发的平台 & 旧互联网平台形态。 \\
产销平台 & 生产和消费合一、按用户意图生成内容的平台 & AI 时代可能出现的新形态。 \\
注意力货币 & 平台以播放、停留、点击等注意力计价 & 传统内容平台的核心货币。 \\
信任货币 & 用户因信任作者/IP/频道而付费或持续消费 & AI 生成后可能更稀缺的价值。 \\
广义 AGI & 宏观上通向通用智能的长期叙事 & 节目中用来区分具体产品机会。 \\
狭义 AGI & 某一类任务或场景中足够通用的智能系统 & 更接近应用公司可验证路径。 \\
\bottomrule
\end{longtable}

\subsection{本章小结}

这些词共同构成一个平台经济框架：数据决定智能边界，context 释放模型能力，生成系统重写内容生产，recipe 和信任替代单点内容和注意力成为更重要的资产。

